\vspace{-3mm}
\subsection{QM9}
\vspace{-3mm}
\label{subsec:qm9_exp}

\paragraph{Dataset.}
%%%The QM9~\citep{qm9_2, qm9_1} dataset (CC BY-NC SA 4.0 license) consisting of 134k small molecules, and the goal is to predict their quantum properties such as energy.
The QM9 dataset~\citep{qm9_2, qm9_1}  (CC BY-NC SA 4.0 license) consists of 134k small molecules, and the goal is to predict their quantum properties.
%The average number of atoms per molecule is 18, and each atom in a molecule comes from the set of H, C, N, O, F atoms and has 3D coordinates.
%%%We follow the data partition used by Cormorant~\cite{cormorant}, which has 100k, 18k and 13k molecules in training, validation and testing sets.
The data partition we use has 110k, 10k, and 11k molecules in training, validation and testing sets.
We minimize mean absolute error (MAE) between prediction and normalized ground truth. % and report MAE on the testing set.


\vspace{-3.25mm}
\paragraph{Training Details.}
Please refer to Sec.~\ref{appendix:subsec:qm9_training_details} in appendix for details on architecture, hyper-parameters and training time.


\vspace{-3.25mm}
\paragraph{Results.}
%%%We mainly compare with methods trained with the same data partition and summarize the results in Table~\ref{tab:qm9_result}.
We summarize the comparison to previous models in Table~\ref{tab:qm9_result}.
Equiformer achieves overall better results across 12 regression tasks compared to each individual model.
%%%Equiformer achieves the best results on 11 out of 12 tasks among models trained with same data partition.
The comparison to SEGNN, which uses irreps features as Equiformer, demonstrates the effectiveness of combining non-lienar message passing with MLP attention.
Additionally, Equiformer achieves better results for most tasks when compared to other equivariant Transformers, which are SE(3)-Transformer, TorchMD-NET and EQGAT.
This demonstrates a better adaption of Transformers to 3D graphs and the effectiveness of the proposed equivariant graph attention.
We note that for the tasks of $\mu$ and $R^2$, PaiNN and TorchMD-NET use different architectures, which take into account the property of the task.
%\revision{This makes the comparison less clear since we use the same architecture for all tasks.}
In contrast, we use the same architecture for all tasks.
We compare training time in Sec.~\ref{appendix:subsec:qm9_training_time_comparison}.

%%%Besides, the different data partition as denoted by $\dagger$ in Table~\ref{tab:qm9_result} has $10\%$ more molecules in the training set and less data in the testing set, and this can benefit some tasks that are more dependent on data partitions.
%\todo{If there is any justification, can write something like which tasks are more important.}